
CORREGIDO en bitxo.tex

For Pt–Ge clusters, we further examine whether CO binding depends not only on the coordination of the adsorption site but also on the chemical identity of its neighboring atoms. We distinguish the effective contributions of Pt and Ge neighbors through \$ECoN\_{\textbackslash{}rm Pt\}^\{0\}$ and \$ECoN\_{\textbackslash{}rm Ge\}^\{0\}$, respectively, and define \$x\_{\textbackslash{}rm Ge\}^\{0\}$ as the fraction of the initial total coordination arising from Ge neighbors.

For \textbackslash{}ce\{Pt6Ge4\^{der\}} (Table\~\ref\{tab\}), site 5/7 has the largest initial Ge contribution (\$x\_{\textbackslash{}rm Ge\}^\{0\}=0.81\$) and exhibits the weakest CO binding (\$BE\_{\textbackslash{}rm ad\}(\textbackslash{}ce\{CO\})=-1.57\$ eV). In contrast, sites 2 and 1/4, with greater effective Pt coordination and smaller Ge fractions, yield three of the strongest binding energies. This suggests that increasing the Ge contribution generally weakens CO adsorption. However, site 9 exhibits strong binding (\$-2.25\$ eV) despite its relatively high initial Ge fraction (\$x\_{\textbackslash{}rm Ge\}^\{0\}=0.75\$). Upon adsorption, its local environment undergoes substantial reorganization, with \$ECoN\_{\textbackslash{}rm Pt\}$ increasing from 0.73 to 1.16 and \$x\_{\textbackslash{}rm Ge\}$ decreasing from 0.75 to 0.61. By comparison, site 5/7 shows smaller changes (\$ECoN\_{\textbackslash{}rm Pt\}$: 0.53 to 0.66; \$x\_{\textbackslash{}rm Ge\}$: 0.81 to 0.77). Thus, the stronger binding at site 9 is accompanied by a greater shift toward Pt coordination, illustrating the limitations of descriptors based solely on the initial geometry.

The \textbackslash{}ce\{Pt5Ge5\^{der\}} cluster further highlights the importance of both the chemical identity and extent of local coordination. Site 4, coordinated almost exclusively by Ge (\$x\_{\textbackslash{}rm Ge\}^\{0\}\simeq1.00\$), exhibits the weakest binding. More interestingly, sites 6/8 and 7/10 have nearly identical initial Ge fractions (\$x\_{\textbackslash{}rm Ge\}^\{0\}=0.75\$), but site 6/8 has a substantially higher \$ECoN\_{\textbackslash{}rm total\}^\{0\}$ (5.58 versus 2.91) and binds CO 0.42 eV more strongly. This coordination difference persists after adsorption (\$ECoN\_{\textbackslash{}rm total\}^\{\ce\{CO\}}=5.69\$ and 2.95, respectively), while both sites undergo a modest shift toward Pt coordination (\$x\_{\textbackslash{}rm Ge\}^\{\ce\{CO\}}=0.70\$ and 0.68). These results indicate that the Ge fraction alone cannot explain the binding-energy differences, as the overall coordination environment also plays a role. The energetic origin of the 0.42 eV difference is examined through the fragment analysis in Section\~\ref\{subsec\}.

For \textbackslash{}ce\{Pt4Ge6\^{der\}}, all adsorption sites are Ge-rich, but site 1, with the highest initial Ge fraction (\$x\_{\textbackslash{}rm Ge\}^\{0\}=0.98\$), exhibits the weakest binding (\$-1.05\$ eV), compared with \$-1.31\$ eV at site 2 and \$-1.46\$ eV at sites 5 and 7, which have greater effective Pt contributions. Notably, site 1 undergoes substantial reorganization upon adsorption, increasing \$ECoN\_{\textbackslash{}rm Pt\}$ to 1.42 and reducing \$x\_{\textbackslash{}rm Ge\}$ to 0.76, yet remains the weakest-binding site. The optimized carbonyl descriptors for the derivative and GM structures are reported in Supporting Tables\~\ref\{tab\} and\~\ref\{tab\}, respectively.

The GM structures show similar trends, although with some exceptions (Supporting Table\~\ref\{tab\}). For \textbackslash{}ce\{Pt6Ge4\^{GM\}}, sites with lower initial Ge fractions (\$x\_{\textbackslash{}rm Ge\}^\{0\}=0.34\$--0.56) bind CO more strongly (\$-1.95\$ to \$-1.66\$ eV) than the more Ge-rich sites 4 and 5 (\$x\_{\textbackslash{}rm Ge\}^\{0\}=0.81\$ and 0.80; \$-1.21\$ and \$-1.22\$ eV). Similarly, in \textbackslash{}ce\{Pt5Ge5\^{GM\}}, sites 1, 2, 4, and 5 (\$x\_{\textbackslash{}rm Ge\}^\{0\}\simeq0.93\$) converge to the same carbonyl minimum, whereas the almost exclusively Ge-coordinated site 3 (\$x\_{\textbackslash{}rm Ge\}^\{0\}\simeq1.00\$) exhibits the weakest binding. In contrast, \textbackslash{}ce\{Pt4Ge6\^{GM\}} shows no clear correlation: sites 2 and 3 have nearly identical Ge fractions (0.92 and 0.91) but exhibit the strongest and weakest binding energies, respectively (\$-1.51\$ and \$-1.05\$ eV).

Overall, these results suggest that a greater Ge contribution to the local Pt environment generally correlates with weaker CO binding. However, deviations from this trend demonstrate that \textbf{initial local composition and coordination are insufficient to predict CO affinity on their own}. The structural changes accompanying adsorption must also be considered, motivating the energetic decomposition of the binding process discussed in the following section.

We next examine the optimized \textbackslash{}ce\{Pt-C\} bond lengths and the electron density at the bond critical point (\$\rho\_{\textbackslash{}rm BCP\}$) to assess how local bonding descriptors relate to CO binding strength. Figure\~\ref\{fig\} reveals a general tendency toward stronger CO binding at shorter \textbackslash{}ce\{Pt-C\} distances, although substantial differences in binding energy are observed even for carbonyls with nearly identical bond lengths.

This behavior is particularly evident for \textbackslash{}ce\{Pt10\^{der\}CO\}. The triplet minima (black filled squares) exhibit nearly identical \textbackslash{}ce\{Pt-C\} distances (1.83--1.84\~\si\{\angstromunit\}), despite a binding-energy variation of approximately 0.50 eV. A similar trend is observed for the quintet minima (black filled triangles), whose binding energies span 0.64 eV over the same narrow distance range. The GM structures (open symbols) exhibit a broader distribution of \textbackslash{}ce\{Pt-C\} distances, reflecting differences in cluster geometry and spin state. The Pt–Ge carbonyls extend the distance range to approximately 1.96\~\si\{\angstromunit\} and provide further evidence that bond length alone does not determine adsorption strength. For instance, sites 6/8 and 7/10 of \textbackslash{}ce\{Pt5Ge5\^{der\}CO\} have similar \textbackslash{}ce\{Pt-C\} distances (1.90 and 1.89\~\si\{\angstromunit\}), yet their binding energies differ by 0.42 eV (\$-2.08\$ and \$-1.66\$ eV, respectively).

The AIM analysis provides a complementary measure of local \textbackslash{}ce\{Pt-C\} bonding through \$\rho\_{\textbackslash{}rm BCP\}$ (Table\~\ref\{tab\}). For the selected strongest- and weakest-binding sites of \textbackslash{}ce\{Pt10\}, \textbackslash{}ce\{Pt6Ge4\}, and \textbackslash{}ce\{Pt5Ge5\} (6 versus 1/9, 2 versus 5/7, and 6/8 versus 4, respectively), stronger CO binding is accompanied by higher \$\rho\_{\textbackslash{}rm BCP\}$. However, this correlation is not universal. Sites 6/8 and 7/10 of \textbackslash{}ce\{Pt5Ge5\} exhibit nearly identical electron densities (0.173 and 0.175 \$e,a\_0\^{-3\}$), despite their 0.42 eV binding-energy difference. Likewise, site 1 of \textbackslash{}ce\{Pt4Ge6\} has a higher \$\rho\_{\textbackslash{}rm BCP\}$ than site 5/7 (0.181 versus 0.171 \$e,a\_0\^{-3\}$), yet binds CO more weakly.

These results demonstrate that \textbf{neither the \textbackslash{}ce\{Pt-C\} bond length nor \$\rho\_{\textbackslash{}rm BCP\}$ alone provides a reliable measure of the overall CO binding energy}. Although both characterize aspects of the local \textbackslash{}ce\{Pt-C\} bond, the adsorption energy also depends on the energetic consequences of cluster reorganization. To disentangle these contributions, we next analyze the interaction and deformation energies associated with CO adsorption (Section\~\ref\{subsec\}).



To complement the ECoN analysis, we decompose the electronic CO binding energy into the frozen-fragment interaction and the energy changes associated with metal-cluster and CO reorganization. A similar approach was employed by Gasper \textit{et al.} to investigate CO adsorption on Pt clusters, identifying the CO–Pt interaction as the dominant contribution, with a smaller contribution from cluster relaxation.\textbackslash{}cite\{gasper:2017\}

\textbackslash{}subsubsection\{Energy decomposition and \textbackslash{}ce\{Pt10\} comparison\}

We first consider three triplet \textbackslash{}ce\{Pt10\^{der\}CO\} minima obtained from the same bare-cluster structure: the strongest-binding minimum reached from site 6, the weakest-binding minimum reached from sites 1/9, and the minimum reached from sites 8/10. For each optimized carbonyl, the interaction energy is calculated using the isolated metal and CO fragments frozen at their geometries in the complex:

\textbackslash{}begin\{align\}

E\_{\textbackslash{}mathrm\{int\}} \&=

E\_{\textbackslash{}ce\{Pt10CO\}}(\textbackslash{}mathbf R\_{\textbackslash{}ce\{Pt10CO\}})

-E\_{\textbackslash{}ce\{Pt10\}}(\textbackslash{}mathbf R\_{\textbackslash{}ce\{Pt10\}}\^{\textbackslash{}ce\{Pt10CO\}})

-E\_{\textbackslash{}ce\{CO\}}(\textbackslash{}mathbf R\_{\textbackslash{}ce\{CO\}}\^{\textbackslash{}ce\{Pt10CO\}}),\textbackslash{}

E\_{\textbackslash{}mathrm\{def\}} \&=

\textbackslash{}left[

E\_{\textbackslash{}ce\{Pt10\}}(\textbackslash{}mathbf R\_{\textbackslash{}ce\{Pt10\}}\^{\textbackslash{}ce\{Pt10CO\}})

-E\_{\textbackslash{}ce\{Pt10\}}(\textbackslash{}mathbf R\_{\textbackslash{}ce\{Pt10\}}\^{0\})

\textbackslash{}right]

\{}+

\textbackslash{}left[

E\_{\textbackslash{}ce\{CO\}}(\textbackslash{}mathbf R\_{\textbackslash{}ce\{CO\}}\^{\textbackslash{}ce\{Pt10CO\}})

-E\_{\textbackslash{}ce\{CO\}}(\textbackslash{}mathbf R\_{\textbackslash{}ce\{CO\}}\^{0\})

\textbackslash{}right],

\textbackslash{}label\{eq\}

\textbackslash{}end\{align\}

where \$\mathbf R\^{\textbackslash{}ce\{Pt10CO\}}\$ denotes the fragment geometry in the optimized carbonyl, while \$\mathbf R\^0\$ refers to the corresponding isolated reference structure (\textbackslash{}ce\{Pt10\^{der\}} or singlet \textbackslash{}ce\{CO\}). The two contributions to \$E\_{\textbackslash{}rm def\}$ are denoted \$E\_{\textbackslash{}rm reorg\}^\{\rm metal\}$      and \$E\_{\textbackslash{}rm reorg\}^\{\ce\{CO\}}\$, respectively. Thus, \$\Delta E\_{\textbackslash{}rm elec\}=E\_{\textbackslash{}rm int\}+E\_{\textbackslash{}rm def\}$, and the zero-point vibrational correction is subsequently added to obtain \$BE\_{\textbackslash{}rm ad\}(\textbackslash{}ce\{CO\})\$.

Importantly, the derivative clusters are not necessarily the most stable isomers but are deliberately retained in their prescribed geometries to enable systematic comparisons. Consequently, metal reorganization can be energetically favorable (\$E\_{\textbackslash{}rm reorg\}^\{\rm metal\}<0\$) when the metal framework adopted upon CO adsorption is lower in energy than the corresponding derivative reference structure.

For \textbackslash{}ce\{Pt10\^{der\}CO\}, the strongest-binding site 6 exhibits a frozen-fragment interaction of \$-2.71\$ eV and a slightly stabilizing metal-reorganization contribution of \$-0.06\$ eV, yielding \$\Delta E\_{\textbackslash{}rm elec\}=-2.74\$ eV. In contrast, the weakest-binding minimum (sites 1/9) has a less favorable interaction (\$-2.46\$ eV) and a positive metal-reorganization contribution (\$0.20\$ eV), giving \$\Delta E\_{\textbackslash{}rm elec\}=-2.24\$ eV. Thus, the 0.50 eV difference in electronic binding energy arises from both a stronger interaction (0.25 eV) and more favorable metal reorganization (0.26 eV) at site 6. The minimum reached from sites 8/10 exhibits similar binding energies and decomposition terms to site 6 (Table\~\ref\{tab\}).

Since the CO-reorganization contribution is small and nearly constant across the selected carbonyls (0.02–0.03 eV), the following discussion focuses on the competition between frozen-fragment interaction and metal-cluster reorganization.

\textbackslash{}subsubsection\{Interaction and reorganization in Pt--Ge clusters\}

The Pt–Ge derivatives reveal distinct balances between interaction and reorganization contributions depending on cluster composition and adsorption site (Table\~\ref\{tab\}).

For \textbackslash{}ce\{Pt6Ge4\^{der\}}, site 2 binds CO 0.68 eV more strongly than site 5/7, primarily due to its 0.79 eV more favorable frozen-fragment interaction. This advantage is only slightly offset by a 0.08 eV less favorable metal-reorganization contribution. Similarly, in \textbackslash{}ce\{Pt4Ge6\^{der\}}, sites 5/7 bind CO 0.41 eV more strongly than site 1, with both the interaction (\$-1.89\$ versus \$-1.66\$ eV) and metal reorganization (0.34 versus 0.52 eV) favoring sites 5/7. These results indicate that, while the frozen-fragment interaction plays a major role in both compositions, metal reorganization can either reinforce or partially compensate for differences in interaction strength.

A markedly different behavior emerges for \textbackslash{}ce\{Pt5Ge5\^{der\}}, where metal reorganization becomes the dominant factor. Site 6/8 binds CO 0.74 eV more strongly than site 4, despite having a less favorable frozen-fragment interaction (\$-1.78\$ versus \$-1.92\$ eV). Instead, its metal-reorganization energy is strongly stabilizing (\$-0.40\$ eV), compared with the substantial energetic penalty at site 4 (\$0.49\$ eV). The resulting 0.89 eV advantage in metal reorganization outweighs the interaction-energy difference and accounts for the stronger overall adsorption at site 6/8.

The decisive role of metal reorganization is particularly evident when comparing sites 6/8 and 7/10 of \textbackslash{}ce\{Pt5Ge5\^{der\}}.} As discussed in Section\~\ref\{subsec\}, these sites have identical initial Ge fractions (\$x\_{\textbackslash{}rm Ge\}^\{0\}=0.75\$) but markedly different total coordination (\$ECoN\_{\textbackslash{}rm total\}^\{0\}=5.58\$ and 2.91, respectively). Despite its higher coordination, site 6/8 binds CO 0.42 eV more strongly (\$-2.08\$ versus \$-1.66\$ eV). Remarkably, the frozen-fragment interaction actually favors the weaker-binding site 7/10 (\$-1.92\$ versus \$-1.78\$ eV). However, this advantage is more than compensated by the metal-reorganization contribution, which is stabilizing at site 6/8 (\$-0.40\$ eV) but destabilizing at site 7/10 (\$0.17\$ eV). The 0.57 eV difference in metal reorganization reverses the trend predicted from the frozen-fragment interaction, yielding electronic binding energies of \$-2.15\$ and \$-1.73\$ eV for sites 6/8 and 7/10, respectively.

Overall, these results demonstrate that the relative importance of frozen-fragment interaction and metal reorganization varies across Pt–Ge compositions. While the interaction term largely determines the binding-energy differences in {Pt6Ge4\^{der\}}, both contributions are important in {Pt4Ge6\^{der\}}, and metal reorganization dominates the selected comparisons in {Pt5Ge5\^{der\}}. In particular, the comparison between sites 6/8 and 7/10 shows that a more favorable Pt–CO interaction does not necessarily lead to stronger overall adsorption, because the energetic consequences of metal-cluster reorganization can reverse the binding-energy trend. This explains why neither initial coordination nor local Pt–CO bonding descriptors alone can fully account for CO affinity.

Although lower coordination generally favors CO adsorption on Pt nanoparticles,\textbackslash{}cite our results show that this trend is modified by the chemical identity of neighboring atoms and the structural response of the cluster. In particular, the relative contributions of the frozen-fragment interaction and metal-cluster reorganization vary across adsorption sites and Pt–Ge compositions, with the latter occasionally reversing the binding-energy trend expected from the interaction alone. Thus, CO affinity cannot be predicted solely from the initial coordination environment but must also account for the energetic consequences of adsorption-induced structural reorganization.